 \documentclass{resume} % Use the custom resume.cls style

\usepackage[left=0.4 in,top=0.4in,right=0.4 in,bottom=0.4in]{geometry} % Document margins
\newcommand{\tab}[1]{\hspace{.2667\textwidth}\rlap{#1}} 
\newcommand{\itab}[1]{\hspace{0em}\rlap{#1}}
\name{Devin De Silva} % Your name
% You can merge both of these into a single line, if you do not have a website.
\address{+1(518) 286-8534 \\ Troy, NY} 
\address{\href{mailto:desild@rpi.edu}{desild@rpi.edu} \\ \href{https://www.linkedin.com/in/devinyasithdesilva}{LinkedIn} \\ \href{https://devindesilva.github.io/Portfolio-Auto/personal}{Website}\\
\href{https://devindesilva.github.io/Portfolio-Auto/personal}{GitHub}
}

%

\begin{document}

%----------------------------------------------------------------------------------------
%	EDUCATION SECTION
%----------------------------------------------------------------------------------------

\begin{rSection}{Education}

{\bf PhD in Computer Science}, \hfill {Expected 2029}\\
Cumulative GPA: 4.00/4.00\\
Rensselaer Polytechnic Institute, NY, USA \\
Second-year 

{\bf MSc in Computer Science}, \hfill {2024-2025}\\
Cumulative GPA: 4.00/4.00\\
Rensselaer Polytechnic Institute, NY, USA 

{\bf B.Sc. Engineering (Hons.) in Computer Science and Engineering} \hfill {2018 - 2023}\\
Cumulative GPA: 3.76/4.20 (First Class Hons.) \\
Faculty of Engineering,
University of Moratuwa, Sri Lanka 
%Minor in Linguistics \smallskip \\
%Member of Eta Kappa Nu \\
%Member of Upsilon Pi Epsilon \\
\end{rSection}

%----------------------------------------------------------------------------------------
% TECHNICAL STRENGTHS	
%----------------------------------------------------------------------------------------
%\begin{rSection}{SKILLS}

%\begin{tabular}{ @{} >{\bfseries}l @{\hspace{6ex}} l }
%Technical Skills & A, B, C, D
%\\
%Soft Skills & A, B, C, D\\
%XYZ & A, B, C, D\\
%\end{tabular}\\
%\end{rSection}

%----------------------------------------------------------------------------------------
% EXPERIENCE	
%----------------------------------------------------------------------------------------

\begin{rSection}{EXPERIENCE}

\textbf{IBM - Research Intern} \hfill May 2025 - Aug 2025\\
\href{https://research.ibm.com/labs/yorktown-heights}{Thomas J. Watson Research Center - IBM} \hfill \textit{YorkTown-heights, NY}
 \begin{itemize}
    \itemsep -3pt {} 
     \item \textbf{DiagnosticIQ} - A Dataset for Benchmarking Decisions LLMs on Industrial Maintenance from faults to fixes.
     \item Maintainance Recommendation engine
 \end{itemize}

\textbf{ML Engineer} \hfill Jun 2023 - Jul 2024\\
\href{https://irononeailabs.com}{IronOne Technologies} \hfill \textit{Austin, TX} \\
\href{https://drive.google.com/file/d/18g-9RbiW1LLCa3f9J8CWWccwRikGCOnR/view?usp=sharing}{Service Letter}
 \begin{itemize}
    \itemsep -3pt {} 
     \item Led the development of a system to predict credit card defaulting with explanations.
     \item Refining the system for extreme class imbalanced data and developing an efficient mechanism \\ for SHAP calculation for self-paced ensemble
    \item Spearheaded the development of a credit risk score model.
    \item Supervised the development of a car damage detection model and breast cancer prediction prototype using genome data (SNPs)
 \end{itemize}
 
\textbf{ML Engineer Internship} \hfill Dec 2021 - Aug 2022\\
\href{https://www.promiseq.com}{promiseQ GmbH} \hfill \textit{Berlin, Germany}
 \begin{itemize}
    \itemsep -3pt {} 
     \item Development of an efficient image stream-based motion detection system (MTS).
     \item Deployment Object detection models and MTS.
     \item Development of a shaky camera detection system
     \item Automating deployment and testing through MLOps.
 \end{itemize}

\end{rSection} 

%----------------------------------------------------------------------------------------
%	PUBLICATIONS SECTION
%----------------------------------------------------------------------------------------

\begin{rSection}{PUBLICATIONS}
\vspace{-1.25em}
\item \textbf{\href{https://link.springer.com/chapter/10.1007/978-981-99-8184-7_34}{TEZARNet: TEmporal Zero-Shot Activity Recognition Network}} \hfill Jul 2022 - Jun 2023
\item \textbf{\href{https://arxiv.org/abs/2411.10273}{Fill in the blanks: Rethinking Interpretability in vision}} \hfill Oct 2023 - Feb 2024
\item \textbf{\href{https://arxiv.org/abs/2507.00050}{SEZ-HARN: Self-Explainable Zero-shot Human Activity Recognition Network}} \hfill Jul 2022 - Jun 2023
\item \textbf{\href{https://arxiv.org/abs/2302.03873}{Geometric Perception-Based Efficient Text Recognition}} \hfill Jul 2022 - Feb 2023
\end{rSection} 
\newpage
%----------------------------------------------------------------------------------------
%	PROJECTS SECTION
%----------------------------------------------------------------------------------------

\begin{rSection}{PROJECTS}
\vspace{-1.25em}

\item \textbf{DoctorAgent: A Clinical Multi-agent Helper for Patient-level Drug
Recommendation}\hfill \href{https://github.com/DevinDeSilva/DoctorAssist}{Github}

MultiAgent system to assist Doctors in selecting drugs for patient diseases. Where the problem is decomposed into simpler problems and each handled by specialized agents.

\textit{Technology}:- ReAct/ Least To Most Prompting, MultiAgent Systems 

\item \textbf{CacheFusion: Hybrid Semantic-Temporal Similarity
KV Cache Clustering}\hfill \href{https://rpiexchange-my.sharepoint.com/:b:/g/personal/desild_rpi_edu/ESXEKPhdWvtHsWfQy-i_tjcBM3565gBZCgLQKYc3wChb5Q?e=f4u2hR}{Report}
\hspace{2mm}
\href{https://github.com/CDFire/CacheFusion}{Github}

Utilizing a semantic and temporal clustering to improve the lossless compressibility of KVCache. The key idea is to identify similar tokens and group them to improve efficiency of LZ4 and ZSTD algorithms.

\textit{Technology}:- Transformers, PyTorch

\item \textbf{Digitgen: Python library for synthetic regular scene text dataset generation }\hfill \href{https://github.com/ACRA-FL/digitgen}{Github}

This is a Python library for creating synthetic scene text and digits with various noises. It is intended to create synthetic datasets for regular scene text recognition and currently has 30k downloads.

\textit{Technology}:- OpenCV, Python, Github CI/CD
\item \textbf{Smart Breadboard for remote lab experiments }\hfill \href{https://drive.google.com/drive/u/0/folders/1xaSNThewvvsGVQ5pdaWTqzAHVvRtRN91}{Video and Report}

Remote Access Breadboard will assist in electronics laboratory work during lockdowns during the pandemic. {\bf 1st place} at the IEEE {\bf PES} national design competition, selected from the IEE {\bf CAS} student design competition 2020-2021, to represent Sri Lanka in the IEEE region 10 (Australia-Asia-Pacific).

\textit{Technology}:- Lucid Chart, Express 
\item \textbf{Forestrack: Deforestation Tracking System}\hfill \href{https://drive.google.com/drive/folders/1QNPDZ0rpAtmN9D_4cJCMls4uN9GD2Hs8?usp=sharing}{Drive}
\hspace{2mm}
\href{https://github.com/Deforestration-Tracker-ADK/forestrack-core}{Github}

Developing a deforestation tracking system that uses a semantic segmentation model to identify changes in Sri Lanka's vegetation coverage using satellite images.

\textit{Technology}:- python, Tensorflow, Django, PlanetAPI

\item \textbf{Predicting Pneumonia Using Human Chest X-ray}\hfill \href{https://medium.com/@devin.18/predicting-the-probability-of-pneumonia-using-chest-x-rays-and-cnns-67457529cf28}{Article}
\hspace{2mm}
\href{https://github.com/DevinDeSilva/Pneaumonia_prediction_using_chest_xrays}{Github}

This task focuses on predicting the presence of pneumonia and its type (viral pneumonia or bacterial pneumonia) using human chest X-ray scans.

\textit{Technology}:- PyTorch, Python , ResNet

% \item \textbf{Druggable Protein Prediction}\hfill \href{https://github.com/DevinDeSilva/DruggableProtienPrediction/blob/main/docs/Prediction_of_Durggable_Proteins_Report.pdf}{Report}
% \hspace{2mm}
% \href{https://github.com/DevinDeSilva/DruggableProtienPrediction/tree/main}{Github}

% Extracted features from protein strings for protein drugability prediction. Feature analysis and results presented in the report

% \textit{Technology}:- Sklearn, Python, Matplotlib
\end{rSection} 

%----------------------------------------------------------------------------------------
\begin{rSection}{Competitions and Hackathonss} 
\begin{tabular}{ @{} >{\bfseries}l @{\hspace{5ex}} l }
Top Ten &  \href{https://github.com/BoTZ-TND/VIPCup-2022}{IEEE VIPCup'22}\\
Champions & \href{https://drive.google.com/drive/folders/11yZMQ4NeLwWPrjqYd2D_5N9qIicXnBvP}{Data Storm 2021: All island open Data science and analytics hackathon}\\
Champions & \href{https://drive.google.com/drive/folders/11yZMQ4NeLwWPrjqYd2D_5N9qIicXnBvP}{Datathon 2021: Inter-University Data Science Competition}\\
Runners up & \href{https://drive.google.com/drive/folders/1H8TTCbwOZNHId4ebueD5T8CW7Qo06PmI}{Datathon 2023: Inter-University Data Science Competition}
\end{tabular}


\end{rSection}

% \begin{rSection}{Technical Skills} 
% \begin{tabular}{ @{} >{\bfseries}l @{\hspace{5ex}} l }
% Languages & Python, R, Javascript \\
% Key Technologies & Pytorch, Tensorflow, SHAP, and Representation Learning
% \end{tabular}


% \end{rSection}

%----------------------------------------------------------------------------------------
\begin{rSection}{Volunteering, Leadership and Extracurriculars} 
\vspace{-0.5 cm}
\item \textbf{International Coordinator, Nalanda College Astronomical Alumni Society}\hfill 2019-2020\\
Working on awareness projects and presenting our projects on international platforms. Dumbara Sky was presented at the 9th Communicating Astronomy with the Public Conference (CAP) 2021

\href{https://www.linkedin.com/posts/nalanda-college-astronomy_virtualcap2021-capconference-activity-6806453733996818432-5cLm}
{Poster} \hspace{0.5mm} \href{https://youtu.be/n-3uiRxUGaw}
{Youtube Presentation}

\item Presenting our project  \textbf{TrailBlazer21} in Communicating Astronomy with the Public Conference (CAP) 2022 

\href{https://drive.google.com/drive/folders/1av3B87YsL-kQWWhhfy79hljFgM_03Sr8?usp=share_link}{Images}
 \hspace{0.5mm} 
\href{https://drive.google.com/file/d/1jzb0nsaYMbM1C9SldCbKqJIjGZ62iaJs/view?usp=share_link}{Participation Invitation}

\end{rSection} 


\end{document}
